\subsection{Clustering no supervisado}
Se aplicaron dos paradigmas de clustering sobre la representaci\'on PCA que conserv\'o el 95\% de la varianza acumulada. Esta elecci\'on redujo el espacio original de 64 coeficientes MFCC a 11 componentes, con lo cual se mitig\'o el efecto de dimensionalidad alta antes de estimar agrupamientos no supervisados.

Para GMM, se evaluaron valores de $k$ en una grilla discreta y se seleccion\'o el n\'umero de componentes que maximiz\'o el coeficiente de Silhouette. Para DBSCAN, se construy\'o una grilla de $\varepsilon$ a partir de distancias al $k$-\'esimo vecino y se seleccion\'o el radio con mayor Silhouette v\'alida sobre puntos no etiquetados como ruido.

\begin{table}[H]
\centering
\caption{Selecci\'on de modelos de clustering mediante Silhouette.}
\begin{tabular}{lrrrrr}
\hline
M\'etodo & Par\'ametro & Clusters & Silhouette & Ruido & ARI \\
\hline
GMM & k=2 & 2 & 0.184 & 0.000 & 0.026 \\
DBSCAN & eps=4.488, min_samples=10 & 2 & 0.436 & 0.056 & 0.004 \\
\hline
\end{tabular}
\end{table}

El coeficiente de Silhouette fue empleado como criterio geom\'etrico porque compara la cohesi\'on intra-cluster con la separaci\'on inter-cluster mediante $(b-a)/\max(a,b)$. Por tanto, valores m\'as cercanos a uno indicaron particiones con menor solapamiento relativo. El ARI se report\'o solo como diagn\'ostico externo, ya que las etiquetas \texttt{species\_id} no fueron utilizadas para ajustar los modelos de clustering.
